\documentclass[abstract=on]{scrreprt}
%\documentclass[abstract=on]{book}
%https://en.wikibooks.org/wiki/LaTeX/Document_Structure

%---------------- lagt til mal -------------------
\usepackage{amsmath}%math env++
\usepackage{bm}%bold math for vectors
\usepackage{amssymb}%for other math symbols
\usepackage[section]{placeins}%forhindrer floats over sections

%--------------- Tilhørende mal ------------------
\usepackage[utf8]{inputenc}
\usepackage{hyperref}
\usepackage{url}
\usepackage{listings}
\usepackage{enumitem}
\usepackage{setspace}
\usepackage[sort=def, nonumberlist, nopostdot]{glossaries}
\usepackage{natbib}
\usepackage{graphicx}
\usepackage{geometry}
\usepackage{pdflscape}
\usepackage{caption}
\usepackage{multirow}
\usepackage{float}
\usepackage{caption}
\usepackage{subcaption}
\usepackage{tikz}
\usepackage{wrapfig}
\usepackage{chngcntr}
\usepackage{comment}
\usepackage[toc, page]{appendix}
\usepackage{titlesec}

% spacing: how to read {12pt plus 4pt minus 2pt}
%           12pt is what we would like the spacing to be
%           plus 4pt means that TeX can stretch it by at most 4pt
%           minus 2pt means that TeX can shrink it by at most 2pt
%       This is one example of the concept of, 'glue', in TeX
\titlespacing\section{0pt}{10pt plus 4pt minus 2pt}{0pt plus 2pt minus 2pt}
\titlespacing\subsection{0pt}{8pt plus 4pt minus 2pt}{0pt plus 2pt minus 2pt}

\makeatletter
\g@addto@macro{\UrlBreaks}{\UrlOrds}
\makeatother

\renewcommand{\abstractname}{} 
\renewcommand{\glossarysection}[2][]{}
\newcommand{\heart}{\ensuremath\heartsuit}

\lstset{basicstyle=\ttfamily}

% paragraph indentation
\setlength{\parindent}{0em} 

% paragraph spacing
\setlength{\parskip}{1.2em}

% table cell spacing
\setlength{\tabcolsep}{1em}
\renewcommand{\arraystretch}{1}

% list spacing
\setitemize[0]{leftmargin=1.5em, itemsep=-0.5em, topsep=-0.4em}

% normal line spacing
\renewcommand{\baselinestretch}{1.1}

 % Next level after subsection
 \newcommand{\paragraphsection}[1]{\paragraph{#1}\mbox{}\\[0.2cm]}
 \renewcommand*{\glossaryname}{Other}
 
\RedeclareSectionCommand[beforeskip=2\baselineskip, afterskip=1\baselineskip]{chapter}


\makeatletter
\renewcommand{\@chapapp}{}% Not necessary...
\newenvironment{chapquote}[2][2em]
  {\setlength{\@tempdima}{#1}%
   \def\chapquote@author{#2}%
   \parshape 1 \@tempdima \dimexpr\textwidth-2\@tempdima\relax%
   \itshape}
  {\par\normalfont\hfill--\ \chapquote@author\hspace*{\@tempdima}\par\bigskip}
\makeatother

\newcommand{\doubleSignature}[3]{
\vspace{1.5cm}

\noindent
\begin{tabular}{lcl}
    \rule{5cm}{1pt} & \hspace{2cm} & \rule{5cm}{1pt} \\
    #2 & & #3
\end{tabular}
}

\newcommand{\singleSignature}[2]{
\vspace{1.5cm}

\noindent
\begin{tabular}{l}
    \rule{5cm}{1pt} \\
    #2
\end{tabular}
}

\definecolor{dkgreen}{rgb}{0,0.6,0}
\definecolor{gray}{rgb}{0.5,0.5,0.5}
\definecolor{mauve}{rgb}{0.58,0,0.82}

\lstdefinelanguage[mips]{java}%
{morestring=[b]",
morekeywords=[1]{},%
comment=[l]\//%
}[keywords,comments,strings]

\lstset{ %
  basicstyle=\scriptsize,       % the size of the fonts that are used for the code
  basicstyle=\ttfamily,
  numbers=left,                   % where to put the line-numbers
  numberstyle=\tiny\color{gray},  % the style that is used for the line-numbers
  stepnumber=1,                   % the step between two line-numbers. If it's 1, each line 
                                  % will be numbered
  numbersep=5pt,                  % how far the line-numbers are from the code
  backgroundcolor=\color{white},  % choose the background color. You must add \usepackage{color}
  showspaces=false,               % show spaces adding particular underscores
  showstringspaces=false,         % underline spaces within strings
  showtabs=false,                 % show tabs within strings adding particular underscores
  frame=single,                   % adds a frame around the code
  rulecolor=\color{white},        % if not set, the frame-color may be changed on line-breaks within not-black text (e.g. commens (green here))
  tabsize=4,                      % sets default tabsize to 2 spaces
  captionpos=b,                   % sets the caption-position to bottom
  breaklines=true,                % sets automatic line breaking
  breakatwhitespace=false,        % sets if automatic breaks should only happen at whitespace
  title=\lstname,                 % show the filename of files included with \lstinputlisting;
                                  % also try caption instead of title
  keywordstyle=\color{blue},          % keyword style
  commentstyle=\color{dkgreen},       % comment style
  stringstyle=\color{mauve},         % string literal style
  escapeinside={\%*}{*)},            % if you want to add a comment within your code
  morekeywords={*,...}               % if you want to add more keywords to the set
}

\renewcommand{\bibname}{References}

% Edit the meta.tex file to change title, group number and author names 
% Fill in the report title, group number and student names here
\newcommand{\mytitle}{Energy consumption of a GPS}
\newcommand{\myauthor}{Jared Showatatek}

\title{\mytitle}
\author{\myauthor}
\date{\today}



\begin{document}
% Caps roman numbers: \pagenumbering{roman} 
\pagenumbering{roman}
% The title page, edit if you want to customize it 

\begin{titlepage}
\begin{center}

 
% Upper part of the page
~\\[0.5cm]
\includegraphics[height=2.0cm]{Images/NTNU_logo.pdf}\\[1cm]  

% Set the title of the Document between two horizontal lines
\hrule ~\\[0.2cm]
{\huge \bfseries \mytitle}\\[0.4cm]		% print the title of the document
\hrule ~\\[1.5cm]

% Additional Information about the document
\begin{minipage}{1\textwidth}
    \centering
	\large
	By
		\myauthor
\end{minipage}

\vfill
{\large \today}
~\\[0.5cm]

% Bottom of the page

\end{center}
\end{titlepage}

% Main matter - edit corresponding file under content/ to change


\begin{minipage}{\textwidth}
\center\Large{Project Assignment:}\\
\end{minipage}
~\\[2cm]
\begin{minipage}{0.8\textwidth}

	\Large Candidate Name: Jared Showatatek\\ 
	\Large Assignment title: Energy consumption of a GPS\\ 
	\Large Assignment text: Model the energy consumption of a GPS\\
	\Large  Co-supervisor: Kjetil Svarstad\\
	\Large Supervisor: Frank Kraemer\\
	
\end{minipage}
\input{Content/02-Abstract.tex}
%%Preface. Here you can describe the process that led to the writing of this report (i.e., the project work process). Here you can mention why you chose this assignment and specific challenges you met with during your work. You can also thank people for help and support during this project.%%%


\chapter*{Preface}



\newpage
%%%%%%%%%TOC%%%%%%%%%%%%%%%%%%%%%%%%%%%%%%%%%%%%%%%%%%%%%%%%%%%%%%%%%%%%%%%%%%%%%%%%%%%%%
\tableofcontents
\newpage
\listoffigures
\listoftables

\newpage

%%%%%%%%%%%%%%%%%%%%%%%INNLEDNING%%%%%%%%%%%%%%%%%%%%%%%%%%%%%%%%%%%%%%%%%%%%%%%%%%%%%%%%
%%%%%%This is normally Chapter 1 and gives an overview of the assignment and why this work is important. If you have chosen to focus mainly on a part of the assignment text, you may write something about this here and explain why. You also normally give a short description of the structure of the rest of the report towards the end of this chapter.  It is important to indicate which parts that are based on your own work. You may even include a list of your main contributions. %%%%%

\chapter{Introduction}
\pagenumbering{arabic}


%%%%%Introduksjon, GPS i IOT verden osv%%%%%%%

Internet of things (IoT) is the network which consists of objects that are connected to the internet such as sensors, vehicles, actuators and other embedded devices. The objects use the internet to deliver data or to be controlled remotely throughout the Internet. IoT is used in the industry to improve the efficiency of operations, safety and security and give valuable insights in analyzing Big data. The use cases for IoT is also relevant for the general consumer as more of the objects in the household are connected to the Internet. Two examples are: controlling the oven temperature with a mobile app or notifying the hospital if a irregular heart rhythm is noticed by the person's pacemaker. IoT is one of the fastest growing trends in technology today, which is highlighted by Gartners Hype cycle of 2017 shown in figure \ref{fig:Gartner}\\


\begin{minipage}[t]{0.8\textwidth}
    \centering
    \includegraphics[width=0.8\textwidth]{Images/Gartner.PNG}\\
    \captionof{figure}{Gartners Hype Cycle 2017}
    \label{fig:Gartner}
\end{minipage}
\\



\subsection{Problem description}
New challenges has emerged as the development of IoT continues to grow, according to \cite{frank} some of them are:

\begin{enumerate}
    \item The scale in terms of units
    \item The constraints in terms of resources: energy, memory, computation
    \item The non-stationary and heterogeneous environments of things.
    
\end{enumerate}

This project will attempt to solve some of the energy constraint IoT designers face today. The project shall explain and build a model of the energy usage of an IoT device which uses GPS. The model shall be so complex that an IoT application can predict the energy cost of getting a positional fix and determine if an alternative method of getting the positional fix is a better substitute. The analyses of the energy usage should also highlight which parameters that are important to include in the energy model. %%%DEVLEOP THIS
\\
\subsection{Motivation}
\cite{frank} also mentions the opportunities in IoT: 
\begin{enumerate}
    \item Many applications only need to provide an overall picture of a situation.
    \item Nodes can	do energy harvesting and the cloud can support nodes with computation.
    \item There	is much data to learn from.
    
\end{enumerate}
This project is part of a ongoing research project at the Faculty of Information Technology and Electronics at NTNU. The research project is called Autonomous Resource-Constrained Things(ART). The aim of the research project is to develop a method for AI to optimize the IoT infrastructure. The first step in developing this method is to analyze and model the energy consumption of the IoT devices. As part of the ART project Ameen Hussain has proposed and evaluated an energy consumption estimation approach for periodic sensing applications running on the IoT devices
devices\cite{Amen}.As the position of an object is one of the most requested information for IoT applications, the usage of GPS has grown substantially.The GPS is used in embedded systems such as watches, trackers, cellphones and cars. GPS can be one of the most power hungry devices in an IoT system. The power budget for an IoT device is often limited, it is therefore useful to have an energy model of the GPS.\\

The student's motivation for choosing this project was to have some insight and experience regarding the different technologies in IoT. The student thinks the experience will be beneficial later as IoT becomes a dominating trend in the Industry. Another motivation was that the Internet of Things combines multiple fields of technology, and challenges the student to apply all of his knowledge of IT and Electronics.
\newpage

\subsection{Methodology}
WRITE THIS LATER
\subsection{Report Structure}

WRITE THIS LATER 
%%%%%%%%%%%%How does a GPS work and draw the state diagram for the GPS.Describe the different alternatives for measuring power consumption of a GPS%%%%%%%%%%%%%%%%%%%%%%%%%





\section{GPS}
The Global positioning system(GPS) is a space based radio navigation system developed by the Unites States Government, and has been operational since 1995.The system provides both timing and geolocation information to a GPS receiver anywhere on the Earth.  The system is not influenced by the number of receivers and can therefore serve an unlimited amount of users. The GPS system can deliver a position which is accurate within 22 meters horizontally if only one receiver is used. If multiple receivers is used positioning accuracy level of the order of a sub-centimeter to a few meters can be obtained \cite{GPS}. 

\subsection{Overview}
GPS consists three segments: space segment, control segment and the user segment. The space segment is a constellation of 24 satellites that are arranged so that  four to ten satellites is visible anywhere on the earth. Each satellite continuously broadcasts a signal composed of two carriers, two digital codes and a navigation message which contains the coordinates of the satellites as a function of the time. Each space vehicle as atomic clocks to ensure the integrity of the navigation message. The codes and the navigation message gets modulated with the two carrier frequencies.  \\If the distance between three satellites are known, the location of the receiver can be determined by measuring the angles with the respect of the each satellite.  GPS needs an additional satellite to account for the clock offset. Figure \ref{fig:GPS} shows the resection that is used by the satellites to determine the position.\\

\begin{minipage}[t]{0.8\textwidth}
\centering
    \includegraphics[width=0.8\textwidth]{Images/gps.PNG}\\
    \captionsetup{justification=centering}
    \captionof{figure}{Resection used by the satellites to determine the position}
    \label{fig:GPS}
\end{minipage}


 The distance between the satellite and the receiver are needed for using the method of resection.The Pseudorange is used for measuring the distance by assuming that the clock of the receiver and satellite are synchronized. It is called for pseudorange since they are actually not perfectly synchronized. The receiver can produce the digital codes at the same time instant that the satellites transmits them from space.The distance can be calculated by measuring the difference in delay of the transmitted signal from the produced one. The clocks is actually not perfectly synchronized and this is reason for the name pseudorange. \\
 A method for determining the velocity of a receiver is by estimating the Doppler frequency of the received signal. The Doppler shift is caused by the relative motion of the receiver and the SV. The two carrier frequencies L1 and L2 that the satellites transmits are generated at 1.5 MHz and 1,2 MHz respectively.  Each SV transmits the L1 and L2 signal but the code modulation for each is different to avoid interference. The two digital code coarse acquisition code (C/A) and precision code(P) consists of a stream of binary digits. The codes are generated using an algorithm and enables a receiver to distinguish the satellites since each have their own set of codes. The navigation message contains the ephemeris, the almanac, the health of the SV, the clock correction, and the atmospheric data.The Ephemeris is the precise information about the orbital position and clock correction for the specific satelite. The Alamanac is a coarse orbital data of all the satellites in constellation. The Ephemeris is only valid for four hours and is sent every 30  second, while the Almanac is valid for 180 days and is sent every 12.5 minutes. \\
 The control segment of the GPS system consists of a number of world wide positioned tracking station and a master control station located in the United States. The control segment tracks the satellites to predict their location, control the atomic clock and satellite data that is transmitted by the SV. \\ The user segment is all GPS receivers that is used to determine their position.
  
  
  \subsection{Acquisition and Tracking}
  
  The state diagram of a GPS receiver consists of two distinct states, the acquisition and tracking states shown in figure \ref{fig:GPS reciever}. The acquisition state is the first state that the receiver goes to after start up. This is the state where the receiver searches and tries to detect the C/A codes from the satellites. After detecting and receiving satellite data from minimum four satellites the receiver goes to the tracking state.
 \begin{figure}[H]
\centering
\includegraphics[width=16 cm]{Project_Report/Images/gps_basics.PNG}
\caption{The state diagram of a GPS reciever}
\label{fig:GPS reciever}
\end{figure}

After synchronizing the clock, the receiver determines the position and continues to search for other satellites signals. The receiver goes back to acquisition state if it doesn't detect minimum four strong satellite signals.
Most receivers uses two separate engines for acquisition and tracking, the acquisition is usually the more power hungry of the two states.The time to first fix (TTFF) depends on three scenarios:
\begin{itemize}
\item Cold start: This is the same as a factory reset. The receiver doesn't have the last positional fix, the time or any satellite data of the constellation. Standard time to first fix  is about 1 minute.
\item Warm start: Last position,time and the Almanac is valid. The receiver does only have to download the Ephirems from the satellites. Standard estimation of TTFF is around 35s.
\item Hot start: When the previous fix was about one second ago, all the data is valid and the estimated time to next fix is 1 second. 

\end{itemize}


\newpage
\input{Content/07-Measurement-methods.tex}
\section{Measuring platform}

The method of using a shunt resistor was chosen to measure the energy consumption of the GPS. The shunt resistor was chosen due to its simple and cheap design but yet its low error rate compared to the Hall effect IC. The other software estimation techniques and cycle aware techniques were not chosen as they either tries to estimate the energy consumption with poor accuracy or they need non available RTL level information about the design. The shunt resistor threats the system that is getting measured as a black box, which enables the measuring platform to be used with any system.  Figure \ref{fig:deploy} shows the deployment diagram of the measuring platform that is used. 

\begin{figure}[H]
\centering
\includegraphics[height=4.5cm]{Project_Report/Images/deploy.PNG}
\caption{The deployment diagram for the rig that is used for measuring the energy consumption}
\label{fig:deploy}
\end{figure}

The platform consist of a micro controller with an embedded GPS, a digital oscilloscope, a shunt resistor, and a PC. The platform is mounted inside a box on a cycle wagon to make the platform mobile. The following section will highlight each parts task. 

\subsection{LoPy with Pytrack}
The LoPy is a micro controller equip with a LoRa, Wifi and BLE technology \cite{LoPy} and is specifically made for IoT use cases. It uses the Espressif ESP32 chip and the user write the program code in Micropython. LoPy has a low-power feature which makes it only use 25uA while still having some its internal peripheral on. The micro controller has UART,SPI, I2C and up to 24 GPIO.\\ The LoPy is equipped with Pycom's Pytrack, which is an extension shield to the LoPy. Pytrack is equipped with a L76-l GPS reciever from Quectel and a 3 axis 12 bit accelerometer \cite{Pytrack}. L76-l is a low power GNSS receiver with have many advanced power saving modes \cite{L76}. The module is equipped with an ARM7 processor and UART for serial communication. The LoPy controls the receiver through the execution of the program code and communicates with the Pytrack through serial communication. TTFF for the receiver is:
\begin{itemize}
\item Cold start:35s
\item Warm start: 30s
\item Hot start: 1s

\end{itemize}



\subsection{Shunt resistor}
\subsection{PC}
A Python script that uses the SDK from Picoscope is executed on the PC. The script is a modified version of Amen Hussain's energy measure script \cite{Amen}.  The scripts measures the energy consumption and plots the data to an Outlook file. The script is included in the appendix. 


\subsection{PicoScope 640 AD}
PicoScope 6000 from Picotecnology is a 4 channel digital oscilloscope with 5GS/s sampling and 2 GSample buffer memory \cite{Pico}. The oscilloscope is equipped with USB 3.0 and supplied with an SDK that enables the user to write their own software. The PicoScope has advanced trigger possibilities and a bandwidth of 500MHz along with a signal generator. The oscilloscope is used to measure the voltage drop across the shunt resistor.  

\newpage
\chapter{Current Consumption}

\section{Measuring with communication}


\begin{table}[h!]
\begin{center}
 \begin{tabular}{||c c c||} 
 \hline
 Waveform(50ms) & Avg Current(mA) & Power(mW)\\ [0.5ex] 
 \hline\hline
 1 & 158    & 24.96 \\ 
 \hline
 2 & 158    & 24.96 \\
 \hline
 3 & 152.3  & 23.19 \\
 \hline
 4 & 165.8  & 27.48 \\
 \hline
 \rowcolor{red}
 5 & 155.1  & 24.05 \\ 
 \hline
 6 & 161    & 25.92 \\ 
 \hline
 7 & 157.6  & 24.83 \\
 \hline
 8 & 149.6  & 22.38 \\
 \hline
 9 & 150    & 22.5  \\
 \hline
 10 & 143.6 & 20.62 \\ 
 \hline
 11 & 151   & 22.80 \\
  \hline
 12 & 157.3 & 24.74 \\
 \hline
 \rowcolor{red}
 13 & 142.3 & 20.24 \\ 
 \hline
 14 & 149.8 & 22.44 \\ 
 \hline
 15 & 150.2 & 22.56 \\
 \hline
 16 & 151.2 & 22.86 \\
 \hline
 17 & 155.4 & 24.14 \\
 \hline
 18 & 142   & 20.16 \\ 
 \hline
 19 & 146   & 21.31  \\
 \hline
 \rowcolor{red}
 20 & 192.7 & 37.13 \\[1ex]
\end{tabular}
\end{center}
\caption{The table shows}
\label{Table:WIFI_ON}
\end{table}

\section{A simpliefied approach}

\begin{table}[h!]
\begin{center}
 \begin{tabular}{||c c c||} 
 \hline
 Waveform(50ms) & Avg Current(mA) & Power(mw)\\ [0.5ex] 
 \hline\hline
 1 & 138.7  & 19.23 \\ 
 \hline
 2 & 130.2  & 16.95 \\
 \hline
 3 & 74.8   & 5.59  \\
 \hline
 4 & 73.22  & 5.36 \\
 \hline
 \rowcolor{red}
 5 & 72.92  & 5.31 \\ 
 \hline
 6 & 73.45  & 5.39 \\ 
 \hline
 7 & 72.63  & 5.27\\
 \hline
 8 & 73.14  & 5.39 \\
 \hline
 9 & 72.84  & 5.30 \\
 \hline
 10 & 73.83 & 5.45 \\ 
 \hline
 11 & 151   &     \\
  \hline
 12 & 157.3 \\
 \hline
 \rowcolor{red}
 13 & 142.3 \\ 
 \hline
 14 & 149.8 \\ 
 \hline
 15 & 150.2 \\
 \hline
 16 & 151.2 \\
 \hline
 17 & 155.4 \\
 \hline
 18 & 142 \\
 \hline
 19 & 146 \\
 \hline
 \rowcolor{red}
 20 & 192.7 \\ 
 [1ex]
\end{tabular}
\end{center}
\caption{The table shows wifi off}
\label{Table:WIFI_OFF}
\end{table}


\section{An improved approach}





\section{Measuring of deepsleep}






















\newpage
 \chapter{Energy Model}
 To optimally use the GPS receiver, we need an model that can predict the energy consumption of the system. In the first subsection, we will first analyze the data from the measurements. In the next subsection, we will present a parameter exploration for the energy model and finally, we will present the model. 
 \section{Data analysis}
 The first step in making an energy model of the GPS receiver is to analyze the data from the measurements and relating it to the theory. The theory section explains how the receiver operates between two distinct phases: acquisition and tracking. The two phases for a receiver can be modelled in a state diagram. 
 
 
\begin{figure}[H]
\centering
\includegraphics[width=16 cm]{Project_Report/Images/Basic_state_GPS.png}
\caption{The state diagram of a GPS receiver}
\label{fig:GPS reciever}
\end{figure}
The receiver will transition to the acquisition state once the update period is set. It will transition from the acquisition state to tracking state once it has acquired a signal from at least four satellites. 
The trigger functionality from table \ref{Table:wifioff} highlights when the receiver has a positional fix. We know from the theory, that the receiver switches to the tracking phase, once it acquires a positional fix. The GPS receiver is therefore in tracking phase the instant the trigger is set. The trigger functionality can't however, inform about later transitions because the receiver changes phases after it has a positional fix to maintain the signal strength. A table with the data of the current consumption right before and after the trigger is set, is generated in an acquisition and a tracking column respectively. The average of each phase is calculated.  Average value of the acquisition and tracking phase is calculated to 80 mA and 72 mA respectively.

Figure \ref{fig:averageacq} shows the scatter plot of the average current consumption during the Acquisition phase. 
 
\begin{figure}[H]
\centering
\includegraphics[width=15 cm]{Project_Report/Images/AcqusitionData.PNG}
\caption{The average values for acquisition phase in a scatter plot }
\label{fig:averageacq}
\end{figure}

Figure \ref{fig:averagetrack} shows the scatter plot of the average current consumption during the tracking phase.

\begin{figure}[H]
\centering
\includegraphics[width=15 cm]{Project_Report/Images/trackingData.PNG}
\caption{The average values for tracking phase in a scatter plot }
\label{fig:averagetrack}
\end{figure}
 
The scatter plots shows how the current consumption is decreasing in the later samples for both phases. The schematics from figure \ref{fig:LoPy_Schematic} shows that the LoPy is using a LiPo battery for voltage supply. Rechargeable batteries  tend to have a discharge curve, which means that the battery voltage changes during discharge. This affect might be the cause of the deviating voltage values in the later measurements. 

The average values are compared with the features from the datasheet \cite{L76}. The datasheet state that it should be a 7 mA difference between the acquisition and tracking phase. This seems to validate the measurements, which shows a decreasing in current consumption from 80 mA to 72 mA after a positional fix.   
\section{Parameter exploration}
The next step in developing an model is to determine which parameters that the model should be dependent on. The parameters in the model will determine the accuracy and abstraction level of the model. We decided to make a simple model with few parameters because of the limited project time.

Based on the theory, we know that the time used in the acquisition state varies accordingly to the validity of the satellite data and signal strength to the satellite. To develop a simple model, we do an optimistic assumption that the time used in acquisition is only dependent on satellite data, and therefore equal to either 35s, 30s or 1s from the specification\cite{L76}. Validity of satellite data is for this reason a parameter for the model. We also assume that the receiver does not switch between the acquisition state and tracking state after acquiring a fix, even if it's kept in tracking state. 

Another parameter is the enabling of deepsleep. The receiver can either be kept in tracking phase continuously or it can be put to deepsleep between update periods. Deepsleep introduces an overhead which consist of initializing the receiver and acquiring signal strength between each updating period. It may therefore exist a scenario where it is better to have continuously fix in tracking, instead of using deepsleep. 

The state diagram in figure \ref{fig:GPS energymodel} is an extension of the general state diagram from figure \ref{fig:GPS reciever}. The state diagram is extended with a deepsleep state and an internal transition if the receiver is kept in tracking state. Timeouts are used for setting time constraints for acquiring a positional fix. 

\begin{figure}[h]
\centering
\includegraphics[width=15 cm]{Project_Report/Images/Energymodel.png}
\caption{The state diagram which represent the energy model}
\label{fig:GPS energymodel}
\end{figure}

\section{The model}
We can derive a prediction of the energy cost of acquiring a fix, by using the state diagram from \ref{fig:GPS energymodel}. The prediction will highlight which update period that is beneficial for having an optimal energy strategy. 
 
The python script in appendix \ref{Appendix:make_energy_model.py} use the method that is derived in this subsection to predict the energy consumption over a range of update periods and durations. An update period is defined as the time from the microcontroller get a request until the next request. The energy consumption over an update period is:
\begin{equation}
E_{period} = P_{period}*T_{period}
\end{equation}
During an update period, the microcontroller transitions between the states in figure \ref{fig:GPS energymodel}. The total energy consumption $E_{fixperiod}$ can therefore be written as the sum of the energy of each state:

\begin{equation}
E_{period} = P_{Initialize}*T_{Initialize} + P_{Acquisition}*T_{Acquisition} + P_{Tracking}*T_{Tracking} + P_{Deepsleep}*T_{Deepsleep}
\end{equation}
If the microcontroller is kept in tracking, the energy used in initialize, acquisition and deepsleep is 0. We omit the initial energy used in initialize and acquisition for the first fix, because it becomes insignificant over a long duration.  

The time used in deepsleep in an update period is dependent on the overhead of waking the microcontroller from deepsleep and acquiring a fix:
\begin{equation}
T_{Deepsleep} = T_{fixperiod} - T_{Initialize} - T_{Acquisition} - T_{Tracking}
\end{equation}
$T_ {Acquisition}$ depends on the validity of satellite data which is given by the fix\_period. The total energy consumption over a duration t is given by the energy consumption of a fix period multiplied by the number of periods during the total duration t.

\begin{equation}
 E_{total} = E_{fixperiod}* \frac{t}{fixperiod}
\end{equation}
\label{equation: energyequation}

 We can use the data from the measurement chapter and calculate the power consumption for each state. This is shown in table  \ref{Table:data for the energy model}. The power is also plotted in a pie diagram in figure \ref{fig:powerconsumption}
\begin{table}[h!]
\begin{center}
 \begin{tabular}{||c c||} 
 \hline
  State &  Power(W) \\ [0.5ex] 
 \hline\hline
  Initialize & 10.201 \\ 
 \hline
 Acquisition & 6.4 \\
 \hline
 Tracking & 5.184 \\
 \hline
 Deepsleep & 0.01025 \\
 [1ex]
 \hline
\end{tabular}
\end{center}
\caption{Power consumption of each state}
\label{Table:data for the energy model}
\end{table}

\begin{figure}[H]
\centering
\includegraphics[height=7.5cm]{Project_Report/Images/powerconsump.PNG}
\caption{The pie diagram of the power in table \ref{Table:data for the energy model}}
\label{fig:powerconsumption}
\end{figure}

The pie diagram shows how the majority of the power is associated with the initializing of the microcontroller. The second biggest contribution to power consumption is the acquisition state, followed by the tracking state. The power consumption in deepsleep is insignificant compared to the other states, which further encourage the use of deepsleep as a power saving strategy. The energy consumption over 1 year with a range of update periods included in table \ref{Table:data for the energy model}. The table is derived by using equation \ref{equation: energyequation}.

\begin{table}[h!]
\begin{center}
 \begin{tabular}{||c c||} 
 \hline
 Energy Consumption & \\
 \hline
  Update Period & Duration = 1 year \\[0.5ex] 
 \hline\hline
  1s & 163482.624 \\ 
 \hline
  7s & 281972.3862 \\
  \hline
  8s & 246766.2040 \\
  \hline
  9s & 219383.6179 \\
  \hline
  10s & 197477.5489 \\
 \hline
  60s & 33182.0320 \\
 \hline
  1800s & 1418.2320 \\
  \hline
  3600s & 870.5803 \\
  \hline
  14399s & 459.8510 \\
  \hline
  14400s & 865.6552 \\
  \hline
  15551700s & 323.4311 \\
  \hline
  15552000s & 323.4959 \\[1ex]
 \hline
\end{tabular}
\end{center}
\caption{Table displaying the energy consumption over different durations and update periods}
\label{Table:energy}
\end{table}


\section{Energy Analysis}

The pie diagram from \ref{fig:powerconsumption} doesn't identify the energy demand of each state. The energy consumption of a state depends on the time that it's active. We know from equation \ref{equation: energyequation} that it's only the time used in acquisition and deepsleep that varies. All the other times in our model is static. The time used in acquisition is an important factor for energy consumption. As explained in the theory chapter, the receiver can be in of three scenarios: 

\begin{itemize}
    \item Hot start: valid Ephemeris \& valid Almanac. Time used in acquisition: 1 s. This is the scenario when the update period is  [1,14399] seconds.
    \item Warm start: invalid Ephemeris \& valid Almanac. Time used in acquisition: 30 s. This is the scenario when the update period is between [14400,15551700] seconds.
    \item Cold start: invalid Ephemeris \& invalid Almanac. Time used in acquisition: 35 s. This is the scenario when the update period is [15551700, $\infty$ ] seconds. 
\end{itemize}

We suspect that the energy consumption of each state will vary for each scenario, since the time in acquisition state varies. 

\subsection{Valid Ephemeris \& valid Almanac}

Figure \ref{fig:firstcase} shows the energy consumption over a duration of 1 year when the update period is 14399 seconds. Th pie diagram shows the contribution from each state to the energy consumption. 

\begin{figure}[H]
\centering
\includegraphics[height=7.5cm]{Project_Report/Images/Allvalid.PNG}
\caption{The energy consumption over 1 year with an update period of 14399 seconds}
\label{fig:firstcase}
\end{figure}

The GPS receiver uses only 1 second in the acquisition state because of the validity of the satellite data. This makes the initializing state, the main factor for the energy consumption. The receiver spends a small amount of time in deepsleep because of the high update frequency. This makes the deepsleep an insignificant contributor to the energy consumption. 

\subsection{Invalid Ephemeris \& valid Almanac}
Figure \ref{fig:secondcase} shows the energy contribution over 1 year for each state when the update period is 14400 seconds. An update period of [14400,15551700] s will make the Ephemeris invalid.

\begin{figure}[H]
\centering
\includegraphics[height=7.5cm]{Project_Report/Images/invalidvalid.PNG}
\caption{The energy consumption over 1 year with an update period of 14400 seconds}
\label{fig:secondcase}
\end{figure}

Comparing figure \ref{fig:firstcase} with figure \ref{fig:secondcase}, we can see that the major energy consumer isn't the initializes state but the acquisition state. This is because of the added time penalty of waiting 30 s instead of 1 s in acquisition state. The added time penalty comes from time used in acquisition to download a valid Ephemeris. The update frequency is high, which makes the receiver spend a small time in deepsleep. This makes deepsleep an insignificant part of the energy consumption.  

\subsection{Invalid Ephemeris \& invalid Almanac}

The energy consumption of each state when both the Ephemeris and Almanac is invalid is show in figure \ref{fig:thirdcase}.

\begin{figure}[H]
\centering
\includegraphics[height=7.5cm]{Project_Report/Images/invalidinvalid.PNG}
\caption{The energy consumption over 1 year with an update period of 15551700 s}
\label{fig:thirdcase}
\end{figure}

The pie diagram in figure \ref{fig:thirdcase} is different from the figure \ref{fig:firstcase} and figure \ref{fig:secondcase}, in that the major energy consumption is during deepsleep. This is because of the low update frequency which cause the system to spend the majority of its time in deepsleep. 

\newpage
\chapter{Discussion and Conclusion}
Table \ref{Table:data for the energy model} from the previous chapter is plotted in figure \ref{fig:energyconsumption}. The plot shows the energy consumption for various update periods over a duration of 1 year.

\begin{figure}[H]
\centering
\includegraphics[height=9.0cm]{Project_Report/Images/energyconsumption.PNG}
\caption{The plot of the energy model from table \ref{Table:energy}}
\label{fig:energyconsumption}
\end{figure}

The plot shows that the energy consumption increases when the update period is increased from 1 second to 7 seconds. This may not be intuitive as the increased update period will cause the system to spend a longer time in deepsleep. The increased energy consumption is caused by the overhead of waking up the system from deepsleep, initializing it and acquiring a fix in acquisition. From figure \ref{fig:firstcase}, we know that the initialing state is a major contributor to energy consumption. An update period that is bigger than  1 second forces the system to repeat this cycle multiple times instead of staying in the tracking state. An update period between (1,7) seconds is not possible because the overhead uses 6 seconds.

The energy consumption decreases when the update period is increased from 7 seconds. At an update period of 10 s, the energy consumption is still higher then an update period of 1 second. By using the algorithm in \ref{code:limit} we can find the limit where the energy consumption is higher than the energy consumption of 1 second.


\lstset{language=Python}          % Set your language (you can change the language for each code-block optionally)
\begin{lstlisting}[frame=single, caption= Code used for finding the beneficial limit]  % Start your code-block

    temp    =   Energy_consumption[4][1]
    optimal     =   Energy_consumption[4][2]
    fix_o = 14400

    while(temp<optimal):  
        T_sleep = fix_o - T_wake - 30 - T_track
        optimal = (P_sleep*T_sleep + P_wake*T_wake + P_acq*1 + P_track*T_track)*t[4]/fix_o
        print(optimal)
        fix_o = fix_o + 1

    print("Limit for using 1s instead is at :" + str(fix_o) + " seconds" )
\end{lstlisting}
\label{code:limit}


The execution of the program code returns 12 seconds. Which means that the energy consumption bigger when an update period in the range [7,12] seconds is instead 1 second. This means that because of the added overhead, it is always more beneficial to use an update period of 1 second which continuously stays in tracking state instead of using an update period of [7,12] seconds. The initial assumption that deepsleep is always a preferred method is therefore inaccurate.

The energy consumption continues to decrease until the update period is 14399 seconds. The energy consumption starts to increase from 14399 seconds. We know from the previous chapter, that an update period that is longer than 14399 seconds will make the Ephemeris invalid. The pie diagram in figure \ref{fig:secondcase} shows how the acquisition state will be the major energy consumer. The addition of the time penalty in acquisition will cause an increase in energy consumption.  From table \ref{Table:data for the energy model} we see that the difference in energy consumption over a year from an update period of 14399 seconds to 14400 seconds is: $865.65$ J $-$ $459.85= 405.8$ J. This shows that according to the model, it is optimum to use an update period that preserves the validity of the Ephemeris and Almanac by having an update period that is slightly less than 14400, instead of having an update period 14400 or slightly more.  The energy consumption decreases from the update period of 14400 until a certain limit. This limit can be found by using algorithm \ref{code:limit}.  The execution of the program returns 57080 seconds, which means that an update period in the range [14400, 57080] seconds will use more energy than an update period of 14399 seconds.  This means that it is more optimal to use an update period of 14399 than an update period in the range [14400, 57080] seconds.

The plot continues to decrease until the updated period is increased to 15551700 seconds. An update period over 15551700 seconds causes the Almanac to be invalid. This means that the system will spend 35 seconds instead of 30 seconds in acquisition state to download the satellite data. The plot in \ref{fig:energyconsumption} shows that the increased energy consumption is insignificant. This is highlighted by table \ref{Table:data for the energy model}, which shows a difference of less than 1 J between an update period 15551700s and 15552000s. The small difference in energy consumption is because of the small penalty that is added to download a valid Almanac. 

In the introduction chapter, we explained the necessity for an energy model for controlling a GPS optimally. This project has explained how the model in figure \ref{fig:energyconsumption} can be used by an application to determine the optimal energy strategy for requesting a positional fix. For instance, if an application wants a fix every 9 second, it can use the model to predict that it is better to have an update period of 1 second.

The time limit of the project compelled us to make a simple model. The consequence of the simple model is that the assumptions might be too oversimplified. The oversimplified model, gives a oversimplified prediction. The assumption that the receiver will use the listed time in acquisition is certainly not always the case. The time used in acquisition is dependent on signal strength, and not including the environment that the receiver operates in, is therefore a crucial oversimplification of the model. Further work should therefore include a probability of the predicted time in acquisition. 

The state diagram from figure \ref{fig:GPS reciever} should also be extended with states that receiver transitions to if it doesn't acquire a fix within the timeout. A failure of acquiring a fix, might require a different strategy as it could inform the application that the receiver is in a environment where signal strength is low. For example a valley or under a tunnel. 

Future work should also establish energy saving techniques that are specific for a certain state and update period. The comparison between the pie diagrams from figure \ref{fig:firstlimit} shows there is a different energy requirement for the states for different update periods. The initialize state is a major energy factor when satellite data is valid. Specific methods should therefore be used for decreasing the energy consumption when a low update period is used. Contrarily, when a greater update period is used, the acquisition is the main energy consumer. 

This project highlights that there are considerable benefits in using an optimal energy strategy for a GPS despite the simplified model. Especially, if the receiver is in use for a long duration. 



% Bibliography - edit references.bib and use the \cite command in text
\bibliographystyle{plain}
\bibliography{Content/references.bib}
\begin{appendices}
\chapter{C027.cpp}
\begin{lstlisting}
#include "mbed.h"
#include "C027.h"
#include "rtos.h"

int power= 0;
char order= 'O';
int BAUDRATE=9600;
int timer;
int MEASURE_TIME= 1000000000;
C027 devkit;
DigitalOut trigger(D0);

void serialhandler(){
    //devkit.gpsPower(true);
    
    // open the gps serial port
    Serial gps(GPSTXD, GPSRXD);
    gps.baud(BAUDRATE);
    
    // open the PC serial port and (use the same baudrate)
    Serial pc(USBTX, USBRX);
    pc.baud(BAUDRATE);
    //devkit.gpsPower(true);
    while(1){
        if(pc.readable()){
            order = pc.getc();
            if(order=='D' && power){
                devkit.gpsPower(false);
                timer= MEASURE_TIME;
                power = 0;
                }
            else if(order=='E' && !power){
                devkit.gpsPower(true);
                timer= MEASURE_TIME;
                power= 1;
                }
            else if(order=='R'){
                    devkit.gpsReset();
                    timer = 0;
                    }  
        }
    if(gps.readable() && pc.writeable()){
        pc.putc(gps.getc());
        }
    }
}
void measure(){
     while(1){
     trigger.write(1);
     wait_ms(10);
     trigger.write(0); //Trekker utgangen ned
     wait_ms(10);
     timer --;
     while(timer<=0){
         wait_ms(1);
         }
    }
}
int main(){
    devkit.gpsPower(false);
    Thread thread_serial;
    Thread thread_measure;
    
    thread_measure.start(measure);
    thread_serial.start(serialhandler);
   while(1){}
    return 0; }
\end{lstlisting}
\label{Appendix:C027.Cpp}

\end{appendices}

\end{document}
